\section*{2.9 延伸阅读}
\addcontentsline{toc}{section}{2.9 延伸阅读}
\label{sec:2.9}

学习线性代数的资源非常丰富，包括 Strang (2003)、Golan (2007)、Axler (2015) 以及 Liesen and Mehrmann (2015) 的教材。我们在本章引言中也提到了一些在线资源。在本书中我们仅介绍了高斯消元法，但求解线性方程组还有许多其他方法；若需更深入的讨论，请参阅 Stoer and Burlirsch (2002)、Golub and Van Loan (2012) 以及 Horn and Johnson (2013) 的数值线性代数教材。

在本书中，我们将线性代数的主题（例如向量、矩阵、线性无关、基）与向量空间的几何相关主题进行了区分。在第 3 章中，我们将引入 \emph{内积（inner product）} ，内积可诱导出 \emph{范数（norm）} 。这些概念使我们能够定义夹角、长度和距离，我们将利用它们来进行 \emph{正交投影（orthogonal projections）} 。事实证明，投影在许多机器学习算法中都起着关键作用，例如 \emph{线性回归（linear regression）} 和 \emph{主成分分析（principal component analysis）} ，我们将在第 9 章和第 10 章分别对它们进行介绍。
